% Nội dung phần 1: Giới thiệu chung và bản đề xuất phân tích.

\section{Giới thiệu chung và bản đề xuất phân tích}

\subsection{Giới thiệu bộ dữ liệu}

Bộ dữ liệu \textbf{e-shop clothing 2008} ghi nhận toàn bộ hành vi click trên một
cửa hàng quần áo trực tuyến của Ba Lan trong năm 2008. Mỗi dòng dữ liệu là một
lượt click; các click có cùng \texttt{session\_id} tạo thành một phiên duyệt web
và biến \texttt{order} cho biết vị trí của click trong phiên đó. Dữ liệu được
công bố kèm nghiên cứu của {\L}apczy\'nski và Bia{\l}ow{\k a}s
(2013)~\cite{lapczynski2013} về hành vi người dùng trong cửa hàng trực tuyến, và
hiện được lưu trữ tại UCI Machine Learning Repository~\cite{uci_clickstream}.

Nhóm chọn bộ dữ liệu này vì ba lý do. Thứ nhất, quy mô đủ lớn (165.474 click của
24.026 session) để tách tập huấn luyện và tập kiểm tra theo thời gian mà vẫn còn
đủ quan sát. Thứ hai, dữ liệu có cấu trúc tuần tự theo session, cho phép đặt một
câu hỏi thống kê phi tầm thường thay vì chỉ mô tả. Thứ ba, dữ liệu kết hợp cả
biến liên tục (giá) lẫn nhiều biến phân loại (danh mục, màu sắc, vị trí ảnh,
trang hiển thị), phù hợp để so sánh hồi quy tuyến tính, GLM và GAM trên cùng một
bài toán.

\begin{table}[H]
\centering
\caption{Tổng quan bộ dữ liệu}
\label{tab:overview}
\begin{tabular}{ll}
\toprule
\textbf{Chỉ tiêu} & \textbf{Giá trị} \\
\midrule
Số click              & 165.474 \\
Số session            & 24.026 \\
Số quốc gia / mã miền & 47 \\
Số sản phẩm           & 217 \\
Số danh mục           & 4 \\
Khoảng thời gian      & 01/04/2008 -- 13/08/2008 \\
Tỷ lệ Continue        & 85,48\% \\
Tỷ lệ Exit            & 14,52\% \\
\bottomrule
\end{tabular}
\end{table}

Bên cạnh các ưu điểm trên, dữ liệu cũng có giới hạn cần nêu ngay từ đầu vì nó
quyết định phạm vi kết luận: không có thông tin giao dịch, giỏ hàng hay thanh
toán; không có giờ click và thời lượng giữa hai click; và tháng 8 chỉ được thu
thập đến ngày 13, trong đó riêng ngày 13/08 chỉ có 200 click. Vì vậy không thể so
sánh tổng lưu lượng tháng 8 trực tiếp với các tháng đầy đủ.

\subsection{Ý nghĩa các biến và biến mục tiêu}

Bảng~\ref{tab:vardict} tóm tắt ý nghĩa và vai trò của từng biến trong phân tích.
Có ba điểm dễ hiểu sai cần lưu ý: \texttt{location} là vị trí của \emph{bức ảnh}
trong sáu vùng của trang web chứ không phải vị trí địa lý của khách hàng;
\texttt{category} tương ứng cột \texttt{PAGE 1} (danh mục chính) còn
\texttt{page} tương ứng cột \texttt{PAGE} (số trang hiển thị 1--5); biến
\texttt{year} bị loại vì chỉ nhận một giá trị duy nhất nên không có biến thiên.

\begin{table}[H]
\centering
\caption{Từ điển biến và vai trò trong phân tích}
\label{tab:vardict}
\footnotesize
\begin{tabular}{>{\raggedright\arraybackslash}p{0.15\textwidth}
                >{\raggedright\arraybackslash}p{0.35\textwidth}
                >{\raggedright\arraybackslash}p{0.40\textwidth}}
\toprule
\textbf{Biến} & \textbf{Ý nghĩa} & \textbf{Vai trò} \\
\midrule
year, month, day & Ngày lịch; không có giờ hoặc thời lượng & Tạo \texttt{time\_30}; \texttt{year} bị loại vì không biến thiên \\[2pt]
order & Thứ tự click trong session & Predictor, dùng $\log_2(\text{order})$ hoặc hàm trơn \\[2pt]
country & Quốc gia hoặc loại tên miền của IP & Predictor factor; mức hiếm được gộp trên train \\[2pt]
session\_id & Mã session & Nhóm dữ liệu và tính robust SE; không làm predictor \\[2pt]
category & PAGE 1: trousers, skirts, blouses hoặc sale & Predictor factor \\[2pt]
product\_model & PAGE 2: mã của 217 sản phẩm & Không đưa vào mô hình chính do có 217 mức \\[2pt]
colour & Màu sản phẩm & Predictor factor \\[2pt]
location & Vị trí ảnh trong sáu vùng của trang; không phải địa lý & Predictor factor \\[2pt]
photography & Kiểu ảnh en face hoặc profile & Predictor factor \\[2pt]
price & Giá sản phẩm, USD & Predictor liên tục \\[2pt]
price\_above\_avg & 1 = cao hơn trung bình danh mục; 2 = không & Tạo hai nhóm cho kiểm định; không dùng cùng \texttt{price} trong mô hình A/B \\[2pt]
page & Trang hiển thị trong website, từ 1 đến 5 & Predictor factor \\
\bottomrule
\end{tabular}
\end{table}

Bộ dữ liệu gốc không có sẵn biến mục tiêu. Nhóm định nghĩa target thống nhất cho
mọi mô hình xác suất trong báo cáo là biến nhị phân \texttt{Exit}:

\[
Exit_{it} =
\begin{cases}
1, & \text{click hiện tại là click cuối của session;}\\
0, & \text{session còn click tiếp theo.}
\end{cases}
\]

Biến \texttt{session\_length} chỉ được dùng để \emph{tạo} target và tuyệt đối
không được đưa vào tập predictor, vì tại thời điểm click thứ $t$ độ dài cuối cùng
của session là thông tin tương lai. Đại lượng được nghiên cứu vì vậy là một xác
suất nguy cơ rời rạc (discrete-time hazard):

\[
P\{Exit_{it}=1\mid \text{session đã tiếp tục đến click }t\}.
\]

Tại \texttt{order} $=t$, mẫu số chỉ gồm các session còn tồn tại đến click $t$.
Do đó mọi kết quả theo \texttt{order} là quan hệ \emph{có điều kiện} trên nhóm
session còn sống sót, không phải bằng chứng rằng việc tăng số click gây ra thay
đổi nguy cơ kết thúc.

\subsection{Các mục tiêu phân tích}

\begin{enumerate}
  \item \textbf{So sánh hai nhóm dạng A/B.} Kiểm định xem xác suất \texttt{Exit}
  có khác nhau giữa click vào sản phẩm có giá cao hơn trung bình danh mục
  (nhóm B) và click vào sản phẩm không cao hơn trung bình (nhóm A) hay không.

  \item \textbf{Phân tích các yếu tố liên quan đến \texttt{Exit}.} Mô tả quan hệ
  giữa xác suất kết thúc và tiến trình session; đánh giá xem đặc điểm của click
  hiện tại và lịch sử duyệt đã quan sát có bổ sung thông tin dự báo hay không.

  \item \textbf{Xây dựng và so sánh mô hình.} So sánh Linear Probability Model
  (LPM), Logistic GLM và Binomial GAM về khả năng phân biệt, chất lượng xác suất
  dự đoán, calibration và khả năng diễn giải.
\end{enumerate}

Cần nhấn mạnh một điểm về thiết kế: nhóm giá \emph{không} được gán ngẫu nhiên cho
người dùng mà là một thuộc tính sẵn có của sản phẩm. Vì vậy Mục tiêu 1 là một so
sánh hai nhóm \textbf{dạng A/B trên dữ liệu quan sát}, không phải một thí nghiệm
A/B có treatment và control. Mọi kết luận rút ra chỉ phản ánh mối liên hệ, không
phải quan hệ nhân quả.

\subsection{Phương pháp và chiến lược phân tích}

\subsubsection{Chiến lược cho Mục tiêu 1}

Giả thuyết được phát biểu hai phía tại mức ý nghĩa $\alpha = 0{,}05$:
\[
H_0: P(Exit=1\mid A)=P(Exit=1\mid B),
\qquad
H_1: P(Exit=1\mid A)\ne P(Exit=1\mid B).
\]
Các bước thực hiện gồm: (i) tính tỷ lệ \texttt{Exit} thô và khoảng tin cậy
Wilson~\cite{wilson1927} cho từng nhóm; (ii) ước lượng risk difference thô kèm
khoảng tin cậy cluster-bootstrap theo session; (iii) ước lượng adjusted odds
ratio và adjusted risk difference từ mô hình logistic có kiểm soát tiến trình
session, thời gian, country và đặc điểm sản phẩm, với sai số chuẩn
cluster-robust theo \texttt{session\_id}~\cite{zeileis2004, cameron2015};
(iv) kiểm tra balance giữa hai nhóm bằng standardized mean difference;
(v) kiểm tra dị biệt hiệu ứng theo \texttt{category} với hiệu chỉnh đa kiểm định
Holm~\cite{holm1979}; (vi) lặp lại phân tích trên ba giai đoạn thời gian để đánh
giá độ ổn định.

\subsubsection{Chiến lược cho Mục tiêu 2 và Mục tiêu 3}

Ba tầng predictor được so sánh trên cùng một pipeline:

\begin{itemize}
  \item \textbf{Baseline:} $\log_2(\texttt{order}) + \texttt{time\_30} + \texttt{country}$.
  \item \textbf{Current-click:} thêm \texttt{price}, \texttt{category},
  \texttt{colour}, \texttt{location}, \texttt{photography} và \texttt{page}.
  \item \textbf{History-enhanced:} thêm số sản phẩm và số danh mục duy nhất đã
  xem, cờ quay lại sản phẩm cũ, cờ chuyển danh mục và chuyển trang, chênh lệch
  giá so với click trước và so với trung bình lịch sử, tỷ lệ quay lại sản phẩm và
  tỷ lệ chuyển danh mục.
\end{itemize}

Mọi đặc trưng lịch sử tại click $t$ chỉ sử dụng dữ liệu từ click 1 đến $t$, nên
không rò rỉ thông tin tương lai. Biến \texttt{order} được đưa vào dưới dạng
$\log_2$ trong LPM và Logistic GLM, và dưới dạng hàm trơn $s(\cdot)$ trong
GAM~\cite{wood2017}, thay vì loại bỏ các session dài nằm ngoài râu boxplot. LPM
và Logistic GLM dùng cluster-robust covariance theo \texttt{session\_id}; GAM
dùng link logit với ước lượng fREML.

Việc chia dữ liệu được thực hiện theo thời gian chứ không ngẫu nhiên, để mô phỏng
đúng tình huống dự báo trong thực tế (Bảng~\ref{tab:split}).

\begin{table}[H]
\centering
\caption{Chia dữ liệu theo thời gian}
\label{tab:split}
\begin{tabular}{llrr}
\toprule
\textbf{Tập} & \textbf{Thời gian} & \textbf{Số click} & \textbf{Số session} \\
\midrule
Train              & Tháng 4--6 & 116.095 & 17.031 \\
Validation         & Tháng 7    & 35.231  & 5.071 \\
Temporal test      & 01--12/08  & 13.948  & 1.903 \\
Loại khỏi đánh giá & 13/08      & 200     & 21 \\
\bottomrule
\end{tabular}
\end{table}

Ngày 13/08 bị loại khỏi mọi đánh giá vì chỉ có 200 click, nhiều khả năng là ngày
thu thập chưa hoàn tất. Toàn bộ quy tắc tiền xử lý được học trên tập train: các
country có dưới 100 click trong train được gộp thành mức \texttt{Other}, rồi áp
dụng nguyên quy tắc đó cho validation và temporal test. Không có session nào xuất
hiện đồng thời trong hai tập bất kỳ, nên không có rò rỉ dữ liệu theo cụm.

Về đánh giá, do lớp \texttt{Exit} chiếm thiểu số (14,52\%), Accuracy không được
dùng làm tiêu chí chọn mô hình. Thay vào đó nhóm dùng ROC--AUC, Average
Precision, log loss, Brier score, Brier Skill Score và hai chỉ số calibration
(intercept và slope)~\cite{steyerberg2010}. Threshold tối đa hóa $F_1$ được chọn
từ prediction out-of-fold của five-fold cross-validation \emph{theo session} trên
tập train, và bị khóa lại trước khi chạm vào validation hay temporal test.
Tháng 7 được dùng để so sánh đặc tả và chọn mô hình; các ngày 01--12/08 chỉ được
mở đúng một lần ở cuối để báo cáo khả năng khái quát hóa.
